\documentclass[11pt, UTF8]{ctexart}
\usepackage{amsmath, amsfonts, amssymb}
\usepackage{extarrows} 
\usepackage{geometry} 
\usepackage{xcolor} 
\usepackage{enumitem} 
\usepackage[most]{tcolorbox} 
\usepackage{fancyhdr} 
\usepackage{diagbox} 
\usepackage{tikz}    
\setlength{\headheight}{13.6pt}
\geometry{a4paper, margin=0.8in} 
\pagestyle{fancy}
\fancyhf{}
\lhead{\color{fblue} \hwdate 作业} 
\rhead{\color{fblue} \today}
\cfoot{\thepage}

\definecolor{fblue}{RGB}{0, 74, 142}
\definecolor{fred}{RGB}{204, 26, 26}

\newtcolorbox{problem}[2]{
    enhanced,
    colback=white,
    colframe=fblue!80,
    arc=2mm,
    boxrule=0.8pt,
    before upper={\textbf{\color{fblue} 问题 #1 (#2)}\hspace{0.5em}},
    before skip=1.5em,
    after skip=1.5em,
    left=3mm, right=3mm, top=2mm, bottom=2mm
}

\newcommand{\hwdate}{5月28日} 
\newcommand{\Prob}{P}
\newcommand{\comp}[1]{\overline{#1}\mskip 3mu}
\newcommand{\ans}[1]{\mathbf{\color{fred}#1}}

\begin{document}

\begin{center}
    {\LARGE \bfseries \color{fblue} \hwdate 作业}
\end{center}

% ---------------- 题目 1 ----------------
\begin{problem}{1}{4. 4}
(1) 设样本 $X_1, X_2, \cdots, X_6$ 来自总体 $N(0,1)$, $Y = (X_1+X_2+X_3)^2 + (X_4+X_5+X_6)^2$, 试确定常数 $C$ 使 $CY$ 服从 $\chi^2$ 分布.

(2) 设样本 $X_1, X_2, \cdots, X_5$ 来自总体 $N(0,1)$, $Y = C(X_1+X_2)/{(X_3^2+X_4^2+X_5^2)^{1/2}}$, 试确定常数 $C$ 使 $Y$ 服从 $t$ 分布.

(3) 已知总体 $X \sim t(n)$, 求证 $X^2 \sim F(1,n)$.
\end{problem}

\paragraph{\color{fblue}解} 
(1) 因 $X_1, X_2, \cdots, X_6$ 是总体 $N(0,1)$ 的样本，故
\[ X_1+X_2+X_3 \sim N(0,3), \quad X_4+X_5+X_6 \sim N(0,3), \]
且两者相互独立. 因此
\[ \frac{X_1+X_2+X_3}{\sqrt{3}} \sim N(0,1), \quad \frac{X_4+X_5+X_6}{\sqrt{3}} \sim N(0,1), \]
且两者相互独立. 按 $\chi^2$ 分布的定义
\[ \frac{(X_1+X_2+X_3)^2}{3} + \frac{(X_4+X_5+X_6)^2}{3} \sim \chi^2(2), \]
即 $\frac{1}{3}Y \sim \chi^2(2)$, 即知 $\ans{C = 1 / 3}$.

(2) 因 $X_1, X_2, \cdots, X_5$ 是总体 $N(0,1)$ 的样本，故 $X_1+X_2 \sim N(0,2)$, 即有
\[ \frac{X_1+X_2}{\sqrt{2}} \sim N(0,1). \]
而
\[ X_3^2+X_4^2+X_5^2 \sim \chi^2(3). \]
且 $({X_1+X_2})/{\sqrt{2}}$ 与 $X_3^2+X_4^2+X_5^2$ 相互独立，于是
\[ \frac{(X_1+X_2)/\sqrt{2}}{\sqrt{(X_3^2+X_4^2+X_5^2)/3}} = \sqrt{\frac{3}{2}} \frac{X_1+X_2}{(X_3^2+X_4^2+X_5^2)^{1/2}} \sim t(3), \]
因此所求的常数 $\ans{C = \sqrt{3 / 2}}$.

(3) 按定义总体 $X \sim t(n)$, 故 $X$ 可表示成
\[ X = \frac{Z}{\sqrt{Y/n}}, \]
其中，$Y \sim \chi^2(n), Z \sim N(0,1)$ 且 $Z$ 与 $Y$ 相互独立，从而
\[ X^2 = \frac{Z^2}{Y/n}. \]
由于 $Z \sim N(0,1), Z^2 \sim \chi^2(1)$, 上式右端分子 $Z^2 \sim \chi^2(1)$, 分母中 $Y \sim \chi^2(n)$, 又由 $Z$ 与 $Y$ 相互独立，知 $Z^2$ 与 $Y$ 相互独立. 按 $F$ 分布的定义得
\[ X^2 \sim F(1,n). \]

\vspace{1.5em}

% ---------------- 题目 2 ----------------
\begin{problem}{2}{4. 6}
设总体 $X \sim b(1,p), X_1, X_2, \cdots, X_n$ 是来自 $X$ 的样本.

(1) 求 $(X_1, X_2, \cdots, X_n)$ 的分布律.

(2) 求 $\sum\limits_{i=1}^n X_i$ 的分布律.

(3) 求 $E(\overline{X}), D(\overline{X}), E(S^2)$.
\end{problem}

\paragraph{\color{fblue}解}
(1) 因 $X_1, X_2, \cdots, X_n$ 相互独立，且有 $X_i \sim b(1,p), i=1,2,\cdots,n$, 即 $X_i$ 具有分布律 $P\{X_i=x_i\} = p^{x_i}(1-p)^{1-x_i}, x_i=0,1$, 因此 $(X_1, X_2, \cdots, X_n)$ 的分布律为
\begin{align*}
P\{X_1=x_1, X_2=x_2, \cdots, X_n=x_n\} &= \prod_{i=1}^n P\{X_i=x_i\} = \prod_{i=1}^n \left[p^{x_i}(1-p)^{1-x_i}\right] \\
&= \ans{p^{\sum_{i=1}^n x_i} (1-p)^{n-\sum_{i=1}^n x_i}}.
\end{align*}

(2) 因 $X_1, X_2, \cdots, X_n$ 相互独立，且有 $X_i \sim b(1,p), i=1,2,\cdots,n$, 故 $\sum\limits_{i=1}^n X_i \sim b(n,p)$, 其分布律为
\[ P\left\{\sum_{i=1}^n X_i = k\right\} = \ans{\binom{n}{k} p^k (1-p)^{n-k}}, \quad k=0,1,2,\cdots,n. \]

(3) 由于总体 $X \sim b(1,p), E(X)=p, D(X)=p(1-p)$, 故有
\[ \ans{E(\overline{X}) = p}, \quad \ans{D(\overline{X}) = \frac{p(1-p)}{n}}, \quad \ans{E(S^2) = D(X) = p(1-p)}. \]

\vspace{1.5em}

% ---------------- 题目 3 ----------------
\begin{problem}{3}{4. 9}
设在总体 $N(\mu,\sigma^2)$ 中抽得一容量为 $16$ 的样本，这里 $\mu,\sigma^2$ 均未知.

(1) 求 $P\{S^2/\sigma^2 \leqslant 2.041\}$, 其中 $S^2$ 为样本方差.

(2) 求 $D(S^2)$.
\end{problem}

\paragraph{\color{fblue}解}
(1) 因为 $(n-1)S^2 / \sigma^2 \sim \chi^2(n-1)$, 现在 $n=16$, 即有 $15S^2/\sigma^2 \sim \chi^2(15)$, 故有
\begin{align*}
p &= P\{S^2/\sigma^2 \leqslant 2.041\} = P\{15S^2/\sigma^2 \leqslant 15 \times 2.041\} \\
  &= P\{15S^2/\sigma^2 \leqslant 30.615\} = 1 - P\{15S^2/\sigma^2 > 30.615\}.
\end{align*}
查 $\chi^2$ 分布表得 $\chi_{0.01}^2(15) = 30.577$, 从而知 $\ans{p = 1 - 0.01 = 0.99}$.

(2) 由 $15S^2/\sigma^2 \sim \chi^2(15)$, 得
\[ D(15S^2/\sigma^2) = 2 \times 15 = 30, \]
即
\[ \frac{15^2}{\sigma^4}D(S^2) = 30, \quad \ans{D(S^2) = \frac{2\sigma^4}{15}}. \]

\end{document}